\documentclass{article}
\usepackage{amsmath, amsthm, amssymb, graphicx}

\title{Proof of the Yang-Mills Gap Theorem}
\author{}
\date{}

\begin{document}

\maketitle

\section*{Chapter 1: Introduction}

\subsection*{1.1 The Yang-Mills Theory}
Yang-Mills theory is a cornerstone of modern quantum field theory, particularly in the description of non-Abelian gauge fields, which are essential in understanding the strong interaction within the Standard Model of particle physics. The theory, formulated by Yang and Mills in 1954, introduces gauge fields that are governed by symmetry groups such as \(SU(2)\) and \(SU(3)\), leading to a rich structure that underlies the fundamental forces of nature.

\subsection*{1.2 The Mass Gap Problem}
One of the most profound challenges in mathematical physics is the Yang-Mills mass gap problem, which is one of the seven Millennium Prize problems posed by the Clay Mathematics Institute. The problem asks whether the quantum field theory based on Yang-Mills fields has a mass gap, meaning that the smallest possible energy of a non-trivial state is strictly positive. This problem is not only fundamental to our understanding of quantum field theory but also has deep implications for the physical spectrum of particles.

\subsection*{1.3 Approach and Objectives}
In this proof, we employ a novel approach to establish the existence of a positive mass gap in Yang-Mills theory. Our strategy involves several key components:
\begin{itemize}
    \item \textbf{Instantons and Energy Quantization:} Instantons, as solutions to the Euclidean Yang-Mills equations, play a crucial role in quantizing energy and preventing massless states.
    \item \textbf{Modified Schwarzschild Metric:} The impact of Yang-Mills fields on black hole metrics is analyzed to show how these fields stabilize black holes and influence entropy and temperature.
    \item \textbf{3s,3t Spacetime and Complex Clocks:} We introduce a novel spacetime framework that supports the existence of a mass gap through phase adjustments during instanton transitions.
    \item \textbf{Morphic Regularization:} A regularization technique that smooths field configurations and ensures stability, further reinforcing the mass gap.
\end{itemize}
The main theorem we aim to prove is the existence of a positive mass gap in non-Abelian Yang-Mills theory.

\subsection*{1.4 Structure of the Proof}
This document is structured into seven chapters:
\begin{itemize}
    \item Chapter 1: Introduction
    \item Chapter 2: Instantons and Tunneling in Yang-Mills Theory
    \item Chapter 3: 3s,3t Spacetime and Complex Clocks
    \item Chapter 4: Modified Schwarzschild Metric with Yang-Mills Fields
    \item Chapter 5: Stability Through Quasi-Normal Modes
    \item Chapter 6: Morphic Regularization
    \item Chapter 7: Conclusion
\end{itemize}

\end{document}